This manuscript addresses the software development challenges for critical reactive systems that
arise from the transition of the automotive industry to Software-Defined Vehicles (\sdv). This
architectural shift replaces distributed networks of specialized \ecus---which communicate over
hardware-centric buses like CAN, designed for small control messages---with centralized computing
platforms. These platforms use high-bandwidth networks, such as Ethernet, capable of managing the
large volumes of data from advanced sensors and applications.
%
This transition, however, introduces significant software challenges. Communication, once handled
efficiently by dedicated hardware, now imposes computational overhead through complex software
protocols. Early experiments at \ampere demonstrated this concretely: attempting to broadcast
messages at fixed intervals---a common practice in legacy systems---quickly saturated the
communication bus. This finding established a critical requirement for \sdv architectures: data must
be transmitted only when its value changes, mandating an event-driven execution model.
%
Concurrently, the industry-wide pursuit of greater cybersecurity and memory safety has led
manufacturers like \ampere to adopt the \Rust programming language. Unlike traditional languages
such as \Ci, \Rust provides compile-time guarantees against entire classes of common
vulnerabilities, including buffer overflows and data races, which are critical for building secure
and reliable automotive systems.
%
However, the adoption of \Rust introduces a new challenge: its low-level nature and steep learning
curve create a significant gap for control engineers, who are accustomed to high-level, model-based
design tools like \Simulink.

To bridge this gap, we introduce \GRust, a domain-specific language designed to reconcile
traditional modeling practices with modern software development needs.
%
\GRust achieves this by combining the strengths of two paradigms. It grounds its design in a
\emph{synchronous dataflow} model~\cite{%
    DBLP:conf/ifip/Kahn74,%
    DBLP:conf/popl/CaspiPHP87,%
    DBLP:conf/tase/ColacoPP17%
}, which provides the deterministic, step-based semantics familiar to control engineers. This
foundation is then enriched with concepts from \emph{functional reactive programming} (\frp)~\cite{%
    DBLP:conf/icfp/ElliottH97,%
    DBLP:conf/haskell/Elliott09,%
    DBLP:conf/padl/WanTH02%
}. The \frp constructs offer a declarative and compositional way to express the asynchronous,
event-driven behavior required by \sdv architectures. This synthesis realizes a \emph{Globally
    Asynchronous, Locally Synchronous} (GALS) architecture~\cite{DBLP:phd/us/Chapiro85}, resulting
in a familiar yet powerful paradigm for system modeling.
%
Furthermore, verifying the safety of high-integrity systems is a critical requirement. The \GRust
ecosystem addresses this through \GReusot, a tool that extends the formal verifier
\Creusot~\cite{DBLP:conf/icfem/DenisJM22}. While \Creusot operates on the generated \Rust code,
\GReusot enables developers to specify and prove safety properties directly within their high-level
\GRust models. This approach ensures that formal verification remains aligned with the engineer's
design intent, preserving the abstraction benefits of the modeling language.

This thesis provides a comprehensive and formal presentation of \GRust, from its design and
theoretical foundations to its practical validation. The primary contributions are as follows.
%
First, we present the design of the \GRust language and its prototype compiler. The compiler
translates high-level, model-based specifications into safe and asynchronous \Rust code that avoids
dynamic memory allocation, ensuring predictable resource usage and bridging the gap between control
engineering and modern software development practices.
%
Second, we establish the theoretical foundations of the language through formal operational and
denotational semantics. These models enable proofs of key properties essential for critical systems:
\emph{productivity}, \emph{determinism}, and the \emph{equivalence} between the two semantic
definitions.
%
Third, we resolve a fundamental conflict between the event-driven execution of services and the
synchronous semantics of components. This conflict leads to a problem we term \emph{oversampling}: a
component can be re-evaluated due to an external trigger, even when its own declared inputs have not
changed. Such re-evaluations can break determinism by causing unintended state updates or producing
spurious outputs.
%
To address this, we introduce a novel \emph{reactive type system}. It statically distinguishes
between two kinds of flows: \emph{continuous flows}, which are inherently tied to the execution
step, and \emph{discrete flows}, which are driven by explicit input changes. The type system ensures
that stateful computations and event-driven logic inside \texttt{when} clauses depend only on
discrete flows. This mechanism acts as a discretization boundary, making component logic resilient
to oversampling. We formalize this guarantee with the \emph{oversampling theorem}. It proves that
for any component typed as \emph{reactive}, an oversampling event does not alter its internal memory
or discrete outputs, thus formally validating the safety and compositional soundness of the
event-driven model.
%
Finally, we demonstrate the practicality of the approach through performance benchmarks. By
evaluating \GRust on representative academic and industrial use cases, we validate its suitability
for real-world deployment in \sdv.

\section{Limitations}
\label{sec:conclusion:limitations}

While this thesis establishes the theoretical foundations and practical viability of \GRust, it is
important to acknowledge the boundaries of the current work. These limitations, which directly
inform the future research directions, are as follows.

First, the prototype compiler, while functional, does not yet implement all features of the language
design. Most notably, the \emph{reactive type system}, which is formally defined, has not been
integrated into the compiler. Consequently, the crucial compile-time guarantees against oversampling
are not yet enforced in practice.

Second, the formal verification of the language semantics was conducted by hand. Although this
approach is sufficient to establish the core properties of the language, the proofs have not been
machine-checked in a formal proof assistant like \Rocq~\cite{rocq} or
Isabelle~\cite{DBLP:books/sp/NipkowPW02}. Such a formalization provides a higher degree of
confidence but was outside the scope of this project.

Finally, the verification capabilities of \GReusot face a key limitation in the compositional
reasoning of stateful components. While properties of a component's own state can be verified,
it is not possible to prove a \texttt{requires} clause of a sub-component if that clause depends on
its internal state. Because the state of sub-components is treated as opaque, their invariants
cannot be used in proofs, which restricts the verification of complex, hierarchical models. This
directly motivates the future work on introducing a public \texttt{invariant} clause.

These limitations clearly define the scope of the contributions and set the stage for the
subsequent development required to transition \GRust from a research prototype to a
production-ready toolchain.

\section{Perspectives}
\label{sec:conclusion:perspectives}

The findings of this thesis suggest that a domain-specific language for modeling offers an effective
path to producing high-quality software for \sdv. The alternative, that is
implementing complex asynchronous logic by hand directly in \Rust, presents considerable obstacles.
Although \Rust provides critical safety guarantees, its low-level nature and the demands of
asynchronous programming require expert-level skills and impose a significant cognitive load on
developers. This approach necessitates a substantial time investment, not only for initial
implementation but also for the extensive testing and qualification phases required in automotive
software. The resulting complexity can compromise code quality and developer productivity, leading
to software that is difficult to maintain and evolve.

A \dsl such as \GRust addresses these challenges by raising the level of abstraction. It allows
engineers to focus on the functional behavior of the system, rather than on the intricate details of
asynchronous programming and memory management in \Rust. As benchmarks on existing \ampere code
suggest (\Cref{sec:grust:benchmarks:grust_vs_rust}), this shift can lead to significant productivity
gains. Furthermore, by making the system design explicit and readable, \GRust enhances long-term
maintainability. The model itself serves as a form of documentation, which facilitates review,
debugging, and modification.
%
This perspective is further strengthened by the promising results from the \GRust prototype, which
exhibits minimal overhead compared to manually implemented code at \ampere
(\Cref{sec:grust:benchmarks:grust_vs_rust}). The prototype compiler was developed with a focus on
demonstrating the practical usability of the language, not on performance optimization. The
efficiency of this unoptimized prototype indicates that the abstraction layer does not impose a
significant performance penalty. This outcome leaves ample room for future optimizations, should
they be required for particularly demanding applications at \ampere.

While this thesis is rooted in the automotive domain, its approach is shaped by the unique
challenges emerging from the architectural shift in the industry. In the transition to \sdv
architectures, the automotive industry is consolidating functions with different safety levels onto
shared hardware, losing the physical isolation of separate \ecus. This trend introduces a
significant risk: a fault in a non-critical component could corrupt a safety-critical one.
The compile-time memory safety guarantees of \Rust are a key enabler in this context, providing
strong software-based isolation to prevent such interference.
%
This specific automotive need makes \Rust adoption compelling, but the \GRust modeling approach has
broader applications. Other domains, such as robotics, have a long history of asynchronous
programming with frameworks like the Robot Operating System (ROS). In ROS, developers manually
implement event propagation by writing nodes, publishers, and subscribers. \GRust offers a
fundamentally different method: it automates the generation of this complex and often error-prone
code from a high-level, declarative model. This suggests that a \GRust-like methodology could
significantly simplify development and improve the maintainability of complex robotic systems.

Ultimately, the choice of a programming paradigm is a strategic decision. Investing in the
development and adoption of a high-level modeling language like \GRust, complemented by a robust
qualification toolchain as discussed in the following section, represents a forward-thinking
strategy. It provides a scalable and sustainable methodology for managing the increasing complexity
of modern automotive software, ensuring that safety, reliability, and maintainability remain central
to the development process.

\section{Future Work}
\label{sec:conclusion:future_work}

The limitations identified in the previous section naturally define a roadmap for future work.
Building on the foundation established in this thesis, the following steps are designed to address
these gaps and enhance the capabilities, safety assurances, and industrial applicability of the
language, transforming the \GRust prototype into a production-ready toolchain for \sdv.

\subsection{Language and Compiler Evolution}

The first priority is to ensure that the practical implementation of the \GRust compiler fully
aligns with its formal specification.

\paragraph{Implementation of the Reactive Type System}
A primary objective for future work is the full implementation of the \emph{reactive type system}
detailed in \Cref{sec:semantics:grsync:react_typing,sec:semantics:grfrp:react_typing}. While the
formal design is complete, the current prototype compiler does not yet enforce its rules. This type
system is critical for guaranteeing the soundness of the execution model at the heart of \GRust. Its
central purpose is to statically distinguish between \emph{continuous} flows (those dependent on
execution time, like \timeop) and \emph{discrete} flows (those only driven by input changes). By
implementing these typing rules, the compiler could formally prove that a component is
\emph{reactive}---that is, resilient to oversampling and safe for use in an event-driven context.
Integrating this system is the necessary step to provide the full safety and predictability
guarantees promised by the \GRust language design.

\subsection{Verification, Validation, and Safety Assurance}

To meet the stringent demands of safety-critical systems, future work must focus on strengthening
the formal verification and validation capabilities of the entire \GRust ecosystem, from high-level
models to the underlying runtime.

\paragraph{Enhancements to \GReusot}
The \GReusot verification wrapper provides a crucial link between the high-level \GRust model and
low-level formal proofs. Future work must address its most significant current limitation: the
verification of stateful components. This would involve introducing a compositional
\texttt{invariant} clause for components. The compiler would then generate inductive proof
obligations to ensure the invariant is established at initialization and preserved at every
subsequent step. To improve usability, support for local assertions within a component body should
also be added. This would allow engineers to guide the prover with intermediate lemmas without
cluttering the public interface. Together, these enhancements would enable the formal verification
of complex, hierarchical models over time, a critical capability for certifying high-\Asil systems.

\paragraph{Integrated Testing and Validation Framework}
To ensure functional correctness and facilitate regression testing, a dedicated testing framework
for \GRust is essential. Inspired by the native test suite of \Rust, this framework would allow
developers to write and run unit tests for components and services directly within the modeling
environment. Engineers could define test cases by providing input streams and asserting expected
output streams, enabling validation at the abstract model level. The compiler would be responsible
for generating and executing these tests, creating a seamless workflow for validating functional
requirements and simplifying the creation of a comprehensive test suite.

\paragraph{Formal Verification of an Asynchronous Runtime for \Rust}
The formal guarantees of productivity and determinism established in this work apply at the level of
individual services. To provide end-to-end assurance, these properties must be extended to the
system as a whole. This requires formally verifying the asynchronous runtime responsible for
orchestrating the execution of \GRust services. The proof must demonstrate that the runtime
correctly preserves the semantics of individual services, ensuring that no data is lost in task
queues and that all temporal constraints are met. A verified runtime would guarantee that the global
system behavior remains predictable and free from concurrency hazards like deadlocks, thus
completing the chain of formal guarantees from model to execution.

\paragraph{Worst-Case Execution Time (WCET) Analysis Integration}
For safety-critical systems, meeting hard real-time deadlines is essential. Future work must
therefore focus on integrating Worst-Case Execution Time (\wcet) analysis into the \GRust toolchain.
By statically analyzing the compiler output---either the generated \Rust code or the underlying LLVM
Intermediate Representation---such tools can compute a sound upper bound for the execution time of a
service reaction. This analysis would provide the formal evidence needed to verify that the system
can meet its timing constraints, a foundational requirement for \iso compliance.

\subsection{Industrialization and Ecosystem Development}

Transitioning \GRust from a research prototype to an industrial asset requires a focus on ecosystem
development and practical validation against another class of real-world demands.

\paragraph{Toolchain for ISO 26262 Qualification}
Supporting industrial adoption requires a toolchain that actively assists with \iso certification.
Future work should enhance the compiler to generate a suite of supporting artifacts. This includes
automatically producing traceability matrices that link requirements to model elements and then to
the generated code. Furthermore, the compiler could parse documentation comments within the \GRust
code to auto-generate design documents. These features would reduce the manual burden of
documentation and provide the verifiable evidence required for safety audits and compiler
qualification.

\paragraph{Benchmarking on Complex Use Cases}
To validate the performance and scalability of \GRust for industrial applications, future work must
focus on benchmarking more complex and computationally intensive systems. Implementing a
representative Advanced Driver-Assistance System (ADAS) would serve as a crucial validation case.
Performance analysis should extend beyond CPU usage to include end-to-end latency, memory footprint,
and communication bus load. Such benchmarks are essential to quantitatively demonstrate the
advantages of the event-driven code generated by \GRust in resource-intensive scenarios, especially
when compared to the traditional, periodic execution models.

\section{Concluding Remarks}
\label{sec:conclusion:remarks}

In summary, this thesis has argued for, and demonstrated the feasibility of, a new approach to
developing software for \sdv. By creating \GRust, we have shown that it is possible to build upon
the well-established foundation of synchronous modeling, which is familiar to control engineers, by
unifying its deterministic rigor with the expressive power of reactive programming.
%
This synthesis provides a bridge from established design principles to the modern demands of safe,
asynchronous, and maintainable systems. By embedding the language directly within the \Rust
toolchain, \GRust makes this bridge a practical reality, enabling a direct translation from
high-level models to safe and performant native code.
%
The formal foundations and practical validation presented here offer a concrete step toward a future
where the complexity of automotive software is managed not by adding layers of manual effort, but by
raising the level of abstraction. This work champions a model-driven methodology, essential for
building the next generation of safe and reliable critical systems.
